\section*{Guía rápida}
\addcontentsline{toc}{section}{Guía rápida}
KalaScript es mi escritura para tomar apuntes: español abreviado, algunos
símbolos matemáticos, veintinueve signos que dicen \emph{cómo se relacionan
las ideas} y veintinueve raíces que se combinan para decir \emph{de qué} hablo.

\subsection*{Por dónde empezar}
\begin{enumerate}
\item Leer la idea y los principios (sección~\ref{sec:idea}).
\item Aprender solo el \textbf{núcleo} en las dos primeras semanas: trece signos
 y, al final, unas pocas raíces (sección~\ref{sec:niveles}). Lo demás espera.
\item Leer las reglas, el estilo telegráfico y, sobre todo, la tabla de
 distinciones (sección~\ref{sec:distinciones}).
\item Imprimir la hoja de consulta (\path{compilado/atlas.pdf}) y tenerla junto
 al cuaderno las primeras semanas.
\item Practicar con \path{compilado/practica.pdf}, usarlo en un curso real y
 medir con las pruebas de la sección~15 del cuaderno.
\end{enumerate}

\subsection*{Archivos}
\begin{description}[style=standard,itemsep=2pt]
\item[\texttt{KalaScript.tex}] Este manual.
\item[\texttt{atlas.tex}] Hoja de consulta de dos páginas para imprimir.
\item[\texttt{practica.tex}] Cuaderno de práctica: copia, lectura, escritura y pruebas.
\item[\texttt{apuntes.tex}] Diez clases de ejemplo en hoja cuadriculada.
\item[\texttt{preambulo.tex}] Configuración y comandos compartidos.
\item[\texttt{simbolos/}] Dibujo de los signos activos, retirados y su trazado.
\item[\texttt{docs/}] Capítulos del manual.
\item[\texttt{anexos/}] Signos retirados e historia de los cambios.
\end{description}
Para compilar todo basta con \texttt{make}; los PDF quedan en \path{compilado/}.
